\uuid{nhIn}
\exo7id{4681}
\titre{exo7 4681}
\auteur{quercia}
\organisation{exo7}
\datecreate{2010-03-16}
\isIndication{false}
\isCorrection{true}
\chapitre{Suite}
\sousChapitre{Convergence}
\cpge{mpsi-mp2i.nombres-reels-suites.h, pcsi.nombres-reels-suites.f, ptsi.nombres-reels-suites.f, mpsi-mp2i.nombres-reels-suites.d, pcsi.nombres-reels-suites.c, ptsi.nombres-reels-suites.c, mpsi-mp2i.nombres-complexes.b, pcsi.nombres-complexes.b, ptsi.nombres-complexes.b, bcpst1.suites-reelles}
\module{Analyse}
\niveau{L1}
\difficulte{}

\contenu{
\question{
Soit $a \in \C\setminus \mathbb{U}$. {\'E}tudier la suite de terme g{\'e}n{\'e}ral :
$u_n = \frac {a^n}{(1+a)(1+a^2)\dots(1+a^n)}$.
}
\reponse{
Si $|a| < 1$, {\`a} partir d'un certain rang, $\left|\frac a{1+a^n} \right|<\alpha < 1$
$ \Rightarrow  \lim = 0$.

Si $|a| >1$, $\frac{u_{n+1}}{u_n} \to 0  \Rightarrow  u_n \to 0$.
}
}
